\unnumberedChapter{Base Station Firmware}

With the base station hardware in place, let’s turn our attention to the firmware.

First of all, the base station is just another LCS node. It builds on the LCS Runtime Library and the Controller Dependent Code library. On top of these foundations, we find locomotive session management, DCC track power management, DCC signal generation, and finally the RailCom communication interface. As an LCS node, the base station exposes attributes and functions through its node and port interfaces, allowing other nodes and applications to control its operation.

The base station is a rather complex node. In particular, DCC signal generation requires the firmware to work closely with the hardware capabilities of the controller chip. We will describe the individual components in the following sections.

From an overall firmware perspective, the base station performs its work in a function that is called periodically by the LCS core library. What sets it apart from other nodes is that the most time-critical task, generating the DCC signal, is interrupt-driven. A central timer generates interrupts whenever the DCC signal needs to change state. It is, in many ways, the heartbeat of the base station. The same principle applies to the RailCom UART input. Everything that belongs to these time-critical communication paths is handled asynchronously through interrupts. All other tasks take place outside the interrupt routines and are comparatively straightforward.

The base station firmware is divided into several key software modules. The following sections provide a high-level overview of these modules and discuss some of the design decisions behind them. We will not attempt to describe every implementation detail here. The source files contain extensive comments documenting the individual components and their operation. Readers interested in the finer details are encouraged to consult the source code alongside this chapter.

\unnumberedSection{Locomotive Session Management}

The base station maintains a list of all active locomotives, as well as a dictionary containing the locomotives known to the system. Before a locomotive or consist can be controlled, it must be found in the dictionary and added to the active session list.

LCS DCC commands, such as setting speed, direction, or functions, are translated into the corresponding DCC packets according to the DCC standards. In addition, the NMRA DCC standards require the periodic refreshing of this information. The session management code simply walks through the active sessions and issues the necessary DCC commands.

Another important part of the base station functionality is the ability to configure DCC decoders. Typically, the locomotive is placed on a separate section of track dedicated to programming. Configuration commands are then used to set the CV values in the decoder. Unlike normal locomotive operation, programming a decoder does not actually require an active locomotive session. However, since programming an engine ultimately comes down to sending DCC packets as well, this functionality is implemented as part of the locomotive session management object.

Locomotive session management interacts closely with the DCC track management objects. Its primary task is to submit the generated DCC packets for transmission on the track. Programming a DCC decoder also requires interaction with the track management, since writing CV values involves detecting changes in the track’s current consumption.

The design deliberately separates DCC packet generation from signal generation. The key interface between the two components is packet submission.
As we will see in the track management section, generating the DCC signal is a typical timer-interrupt-driven task. Whenever a new bit needs to be transmitted, the interrupt handler sets the signal lines that drive the power unit and programs the timer for the next signal transition.

There are designs that let the interrupt handler interact directly with the locomotive session table to retrieve the next bit of the current packet. We have chosen a different approach, with a clear separation between the two responsibilities. Locomotive session management takes care of producing DCC packets and submitting them for transmission. Turning those packets into a stream of bits on the track is entirely the responsibility of the companion DCC track management module.

As the layout grows and the number of running locomotives increases, the available DCC bandwidth may eventually become a limiting factor. There are roughly 100 DCC packets that can be transmitted to the track per second. Using the DCC bus efficiently is therefore important, particularly when the base station has to manage a large number of locomotives.

For a small layout, simply sending out the commands as they arrive and periodically refreshing the active sessions works perfectly well. However, if the base station has to deal with hundreds of locomotives, some form of packet scheduling and priority management may become necessary.

For example, deceleration commands should receive a higher priority than acceleration commands. It is more important to stop a locomotive than to accelerate one. The function refresh option also needs to be considered. Speed and function group refresh packets should be interleaved to ensure that speed commands are repeated frequently enough. A locomotive with a speed setting of zero has a rather low priority, as it is not considered to be actively running.

The DCC standards require each active decoder, that is, a decoder controlling a locomotive with a non-zero speed setting, to receive a refresh at least every 2.5 seconds.
There are also a number of global requirements for the base station. In general, a valid DCC packet must be transmitted at least every 30 milliseconds to indicate that the DCC protocol is still active. Furthermore, there should be a gap of at least 5 milliseconds between two packets addressing the same decoder.

Every five seconds, the base station must also broadcast information about its capabilities. For example, it can announce support for 128-step speed control. There are capability and configuration attributes for stationary decoders as well. Finally, the base station periodically broadcasts the layout time. Decoders can use this information to implement time-dependent features.

For now, our packet scheduling scheme is deliberately simple. An incoming LCS DCC management message directly results in the corresponding DCC command being submitted for transmission. The refresh function simply walks through the active session list and emits a refresh packet for each locomotive or consist.

More experience with really large layouts will be needed before introducing a more sophisticated scheduling scheme. For the time being, the simple approach keeps the implementation straightforward while providing all the functionality we need.

\unnumberedSection{Locomotive Dictionary}

All locomotives need to be known to the base station. They are maintained in a simple dictionary.

Commercial base stations typically record quite a bit of information about each locomotive: its name, a nice picture, its DCC address, and various other attributes. Our base station takes a somewhat different approach. We assume that somewhere in the LCS environment there is a dictionary describing locomotives and other layout-related information in great detail. Our own dictionary only keeps the information needed for locomotive session management.

The most important item is, of course, the DCC address. In addition, we keep the initial DCC function settings and the locomotive name.
Whenever the base station receives a DCC command for a locomotive, it first checks the active session list. If the locomotive is not already active, the dictionary is consulted. If an entry is found, the locomotive is added to the active session list. From that point on, it becomes part of the regular DCC refresh cycle.

The locomotive dictionary is implemented using the LCS node global attribute space. The entries have been carefully crafted to fit in the 16-bit words of the attribute space. Loading the dictionary or adding a locomotive is merely a setting a list of attributes.

\unnumberedSection{Base Station Global Functions}

In addition to locomotive session management and DCC track management, the base station is responsible for broadcasting global information, such as system capabilities and layout time.
All decoders can make use of the layout time. The base station is responsible for transmitting this information periodically.

Furthermore, the base station needs to announce its global capabilities and configuration options. For example, it communicates which DCC standards are supported, whether 128-step speed control is available, and whether programming on the main track is enabled. The base station must broadcast this kind of information every five seconds.

\unnumberedSection{DCC Track Signal Management}

The base station generates the DCC signals for both the MAIN and the PROG track outputs.
To recap, a DCC ONE bit is represented by a square-wave signal with a period of 58 microseconds, while a DCC ZERO bit has a period of 116 microseconds. The signal generation logic could therefore operate in ticks of 58 microseconds.

However, supporting RailCom introduces an additional requirement: the cutout period. A cutout starts with a 29 microsecond ONE pulse, followed by the cutout period itself, and ends with another 29 microsecond pulse to complete the signal transition. Conveniently, this fits perfectly into a timing scheme based on 29 microsecond ticks.

A DCC ONE bit therefore consists of two 29 microsecond ticks with the signal lines in the positive state, followed by two ticks in the negative state. This gives us a complete 58 microsecond positive half-wave followed by a 58 microsecond negative half-wave.
The software implementation in our base station consists of a track signal state machine and a timer that can be programmed to generate interrupts at intervals of multiples of 29 microseconds.

Whenever an interrupt occurs, the interrupt handler first sets the new signal levels and then determines when the next interrupt is required. Since we have two tracks and control both state machines using a single timer, the next interrupt is scheduled for whichever track requires an earlier signal transition.

For example, if one track is scheduled to change its signal level in two ticks and the other in four ticks, the timer is programmed to interrupt again after two ticks. At that point, the first track performs its scheduled transition, while the second track simply decrements its remaining tick count from four to two.
Running both state machines from a single timer also requires splitting the work performed by each state machine into two parts.

First, each state machine takes care of setting the hardware signal levels. It is essential that both tracks receive their signal updates as quickly as possible to maintain accurate DCC timing.

There may, however, be additional work to perform. For example, the state machine might need to retrieve the next bit from the current DCC packet or, once all bits have been transmitted, line up the next packet for transmission. Each state machine can therefore request additional processing. Once both tracks have completed their immediate signal updates, these follow-up actions are handled.

\unnumberedSection{DCC Track Power Management}

In addition to retrieving the next bit or packet, the base station needs to control track power and measure the actual power consumption.

Power consumption measurement essentially involves reading a voltage using the controller’s ADC hardware. The Raspberry Pi Pico has multiple ADC channels, but only one channel can be converted at a time. Nevertheless, a measurement takes only a couple of microseconds, so we can easily measure both tracks within a single interrupt cycle.

There are several ways to perform these measurements. The simplest approach is to sample the current consumption in sync with the DCC signal and store the result. The LCS base station power module does exactly this during a DCC ZERO bit, in the middle of the first half-wave period. The measured value is stored in the corresponding DCC track object and overwritten with each new measurement. Depending on the packet content, a few hundred samples per second can be collected.

While this approach works reasonably well, there is a hidden challenge. The measurement captures the instantaneous current consumption, which depends on the actual power drawn by the decoder. Since decoders typically drive the locomotive motor using PWM, the measurement also depends on when the sample is taken.

With a decoder operating at a relatively low PWM frequency, there is a good chance that a measurement will coincide with the OFF period of the PWM signal. The resulting value would then be considerably lower than the actual average current consumption.
It is therefore important to measure frequently and calculate a mean value across a sufficiently large number of samples. One suitable approach is to calculate the Root Mean Square (RMS) value.

Of course, we also need to make sure that the controller is not overwhelmed by all the processing involved. Our approach is to sample frequently and store the resulting digital values. Short-circuit detection and a simple high-water-mark calculation can be performed with very little overhead.

Only when an RMS value is required do we perform the more expensive calculation involving the sum of the squared samples, followed by division by the number of samples and calculation of the square root. The resulting value is then converted into milliamps.

The RMS calculation uses 32-bit unsigned integer arithmetic on the Pico. This provides sufficient precision for our measurements while avoiding the overhead of floating-point calculations.
Based on these measurements, the track state machine also monitors power consumption for overload conditions. If the current exceeds the configured limit, the track is shut down.

After a short delay, the base station attempts to restore track power. If the restart again results in an overload, the procedure is repeated. After several unsuccessful attempts, the track is shut down permanently and manual intervention is required.

For the programming track, current measurement serves another very important purpose.
When a decoder is programmed using DCC CV access commands, it acknowledges a successful operation by temporarily increasing its current consumption. According to the DCC standards, an increase of 60 mA for a period of 5 to 7 milliseconds constitutes a positive acknowledgement.

Accessing a CV consists of sending a sequence of DCC packets. First, a number of RESET packets are transmitted, followed by the actual CV commands and then another sequence of RESET packets.
During the initial RESET packets, the base station measures the current consumption and establishes a baseline representing the decoder’s normal operating current. When the decoder subsequently acknowledges a request by increasing its current consumption, this baseline is used to detect the acknowledgement.

\unnumberedSection{RailCom Support}

To support RailCom, the power unit is instructed to short-circuit the track for a defined period before transmission of the next DCC packet. This is known as the cutout period.

The DCC decoder in the locomotive detects this cutout period and uses it to transmit a RailCom datagram. During this time, the controller attempts to read the incoming datagram bytes through the serial input line.

Both analog current measurement and RailCom serial input follow a similar implementation pattern: a setup routine initiates the operation, and a corresponding interrupt routine processes the results with as little overhead as possible. After all, maintaining accurate DCC signal timing remains our highest priority.

A RailCom datagram provides two channels. Channel 1 is normally reserved for reporting the locomotive (cab) address. The DCC track management software determines the cab address and stores it in a base station attribute.

Channel 2 provides additional information transmitted by the decoder. Depending on the RailCom application, this can include information such as DCC command acknowledgements or other decoder data. This information is also made available through attributes. In addition, a callback can be registered to allow the application to react to incoming data.

\unnumberedSection{LCS message interface Support}

Being an LCS node, the base station emit LCS messages and reacts to configured events. The message interface module will receive LCS messages and translate them to locomotive session and DCC track commands. 

??? what best to say / show ?

\unnumberedSection{Base Station Serial Commands}

Like any other LCS node, the base station provides a serial command interface. So far, we have encountered this interface mainly as a means of inspecting LCS data structures and issuing simple commands to a node for testing and debugging.

One of the capabilities provided by the LCS library is the ability to invoke a command callback when a command that is not part of the standard LCS command set is received. The base station uses this feature to implement its own command interface.

Naturally, nobody is going to sit in front of a computer and type cryptic commands to control a locomotive. However, serial commands provide an excellent interface for other software, such as the JMRI tools suite or layout management programs like RocRail.

Both of these systems offer interfaces based on the \texttt{DCC++ / DCC-EX} command syntax. Since our base station implements the relevant commands, integration with these systems is relatively straightforward.

The command interface is an important part of the base station functionality and deserves a chapter of its own. We will therefore take a closer look at it in the next chapter.


\unnumberedSection{Base Station Session Management Items}

Being an LCS node, we think in attributes and functions items.

??? keeps also the overarching items ?

\begin{longtable}{@{}|l|l|l|p{0.5\linewidth}|@{}}
    %\caption{Bus Connector Pins}
    \toprule
    \textbf{Item} & \textbf{Arg0} & \textbf{Arg1} & \textbf{Purpose} \\
    \midrule
    \endfirsthead
    \toprule
    \\textbf{Item} & \textbf{Arg0} & \textbf{Arg1} & \textbf{Meaning} \\
    \midrule
    \endhead
    \midrule
    \multicolumn{2}{r}{\textit{Continued on next page}} \\
    \midrule
    \endfoot
    \bottomrule
    \endlastfoot
    - & - & - & - \\
    \midrule
    - & - & - & - \\
\end{longtable}

\unnumberedSection{Base Station DCC Track Management Items}

A DCC track is represented as an LCS port. There are two ports, one for MAIN and one for PROG. Each port 
features a list of attributes.


\begin{longtable}{@{}|l|l|l|p{0.5\linewidth}|@{}}
    %\caption{Bus Connector Pins}
    \toprule
    \textbf{Item} & \textbf{Arg0} & \textbf{Arg1} & \textbf{Purpose} \\
    \midrule
    \endfirsthead
    \toprule
    \\textbf{Item} & \textbf{Arg0} & \textbf{Arg1} & \textbf{Meaning} \\
    \midrule
    \endhead
    \midrule
    \multicolumn{2}{r}{\textit{Continued on next page}} \\
    \midrule
    \endfoot
    \bottomrule
    \endlastfoot
    - & - & - & - \\
    \midrule
    - & - & - & - \\
\end{longtable}


\unnumberedSection{Base Station Locomotive Dictionary Items}

Being an LCS node, we think in attributes and functions items. The dictionary is an array of 
locomotive dictionary entries, each of which is a set of attributes.

\begin{longtable}{@{}|l|l|l|p{0.5\linewidth}|@{}}
    %\caption{Bus Connector Pins}
    \toprule
    \textbf{Item} & \textbf{Arg0} & \textbf{Arg1} & \textbf{Purpose} \\
    \midrule
    \endfirsthead
    \toprule
    \\textbf{Item} & \textbf{Arg0} & \textbf{Arg1} & \textbf{Meaning} \\
    \midrule
    \endhead
    \midrule
    \multicolumn{2}{r}{\textit{Continued on next page}} \\
    \midrule
    \endfoot
    \bottomrule
    \endlastfoot
    - & - & - & - \\
    \midrule
    - & - & - & - \\
\end{longtable}

\unnumberedSection{Summary}

The base station, regardless of how it is implemented, is a key component of any digital model railway layout. It is primarily responsible for locomotive session management and DCC track signal generation. There are many ways to build such a device. Some designs put everything into a nice housing with a display, switches and other controls. Others simply place the LCS node somewhere on the layout, leaving control and management to other devices or software. All these approaches work just fine. And just like any other LCS node, the base station exposes its status and control functions through node and port attributes.

The firmware follows a modular approach, separating locomotive session management, track power management, DCC signal generation and RailCom communication. This separation keeps the individual tasks manageable, even though the base station itself is a rather complex device. It also illustrates how the LCS architecture can accommodate quite sophisticated hardware while maintaining a consistent programming and communication model.

The base station was also our first major LCS node involving a considerable amount of hardware and software. Its development went through several iterations and provided us with quite a few learning experiences. Eventually, after moving to the Raspberry Pi Pico controller and progressing from initial breadboard experiments to dedicated PCBs, we arrived at the base station presented in this chapter.

The lessons learned along the way proved extremely valuable for the overall LCS concept, particularly with respect to hardware architecture and firmware design. There was also quite a bit of reading involved, from controller chip specifications to, perhaps most importantly, the DCC standards themselves.

The appendix contains a collection of useful references and web links. I highly recommend reading the electrical and DCC-related standards published by the RailCommunity organization. They provide an excellent foundation for understanding the technical requirements behind the design decisions discussed in this book.